\title{Επισκόπηση Λειτουργίας Συστήματος Μεταφοράς Ηλεκτρικής Ενέργειας}

\section{Θεωρία}
\label{kef-01-theoria}

Το κεφάλαιο αυτό πραγματεύεται ένα ερώτημα που φαίνεται απλό χωρίς να είναι: με ποιον
τρόπο διατηρείται σε ασφαλή λειτουργία ένα σύστημα ηλεκτρικής ενέργειας. Η απάντηση
συνίσταται σε έναν κλειστό βρόχο τεσσάρων σταδίων — μέτρηση του συστήματος, μετατροπή των
μετρήσεων σε αξιόπιστη εικόνα, αξιολόγηση της ασφάλειας της εικόνας αυτής, και δράση. Οι
ενότητες που ακολουθούν εξετάζουν διαδοχικά τα στάδια αυτά. Το ζητούμενο δεν είναι η
απομνημόνευση των ακρωνυμίων, αλλά η κατανόηση του λόγου ύπαρξης κάθε σταδίου και των
συνεπειών της απουσίας του.

Ως προς την ορολογία, το κεφάλαιο ακολουθεί την επίσημη ελληνική απόδοση των ευρωπαϊκών
κανονισμών για τη λειτουργία του συστήματος και την εξισορρόπηση (Κανονισμοί (ΕΕ) 2017/1485
και 2017/2195), την οποία χρησιμοποιεί και ο ΑΔΜΗΕ. Όπου η ελληνική πανεπιστημιακή
βιβλιογραφία χρησιμοποιεί διαφορετικό όρο, αυτός δίνεται σε παρένθεση κατά την πρώτη
εμφάνιση, μαζί με τον αγγλικό όρο και το καθιερωμένο ακρωνύμιο.

\subsection{Ο ρόλος του Συστήματος Μεταφοράς}

Το Σύστημα Μεταφοράς Ηλεκτρικής Ενέργειας (ΣΜΗΕ) είναι το δίκτυο υπερυψηλής και υψηλής
τάσης που διασυνδέει τους σταθμούς παραγωγής με τα κέντρα κατανάλωσης και με τα γειτονικά
συστήματα. Στην Ελλάδα, το Ελληνικό Σύστημα Μεταφοράς Ηλεκτρικής Ενέργειας (ΕΣΜΗΕ) λειτουργεί στα
επίπεδα των 400 kV, 150 kV και 66 kV. Τη διαχείρισή του ασκεί ο Ανεξάρτητος Διαχειριστής
Μεταφοράς Ηλεκτρικής Ενέργειας (ΑΔΜΗΕ), μέλος του Ευρωπαϊκού Δικτύου Διαχειριστών Συστημάτων
Μεταφοράς Ηλεκτρικής Ενέργειας (ΕΔΔΣΜ ηλεκτρικής ενέργειας, ENTSO-E), ενώ το σύστημα ανήκει
στη συγχρονισμένη περιοχή της Ηπειρωτικής Ευρώπης (Continental Europe, CE).

Ο λόγος για τον οποίο η μεταφορά γίνεται σε τόσο υψηλή τάση προκύπτει άμεσα από τις
θεμελιώδεις σχέσεις. Για δεδομένη μεταφερόμενη ισχύ $P = \sqrt{3}\,V I \cos\varphi$, το
ρεύμα είναι αντιστρόφως ανάλογο της τάσης, ενώ οι απώλειες στους αγωγούς είναι ανάλογες
του $I^2$. Διπλασιασμός της τάσης υποτετραπλασιάζει επομένως τις απώλειες μεταφοράς. Σε
αυτό οφείλεται και η ιεραρχική οργάνωση του συστήματος: υπερυψηλή τάση για τις μεγάλες
αποστάσεις και διαδοχικοί υποβιβασμοί προς την πλευρά του καταναλωτή.

\subsubsection*{Οι τρεις ταυτόχρονες απαιτήσεις}

Ο Διαχειριστής Συστήματος Μεταφοράς (ΔΣΜ· Transmission System Operator, TSO) καλείται να
ικανοποιεί ταυτόχρονα τρεις απαιτήσεις. Η μεταξύ τους ένταση είναι εκείνη που προσδίδει
στο πρόβλημα τη δυσκολία του:

\begin{itemize}
\item \textbf{Ισοζύγιο ισχύος.} Η ηλεκτρική ενέργεια δεν αποθηκεύεται σε σημαντικές
      ποσότητες εντός του συστήματος. Παραγωγή και ζήτηση οφείλουν να εξισώνονται
      \emph{συνεχώς}. Σε αντίθετη περίπτωση, η διαφορά καλύπτεται προσωρινά από την
      κινητική ενέργεια των στρεφόμενων μαζών, με τίμημα τη μεταβολή της συχνότητας: ένα
      σύστημα που καταναλώνει περισσότερο από όσο παράγει επιβραδύνεται κυριολεκτικά.
\item \textbf{Επιχειρησιακή ασφάλεια.} Το σύστημα οφείλει να παραμένει εντός των ορίων
      επιχειρησιακής ασφάλειας (operational security limits) — θερμικών ορίων του
      εξοπλισμού, ορίων τάσης, ορίων ευστάθειας — όχι μόνο στην τρέχουσα κατάσταση, αλλά
      και μετά από κάθε πιθανό μεμονωμένο απρόβλεπτο συμβάν (contingency), όπως η απώλεια
      μιας γραμμής ή μιας μονάδας. Η διάκριση αυτή μεταξύ της τρέχουσας κατάστασης και της
      κατάστασης μετά το συμβάν είναι θεμελιώδης και οδηγεί στο κριτήριο N-1, το οποίο ορίζεται αμέσως παρακάτω.
\item \textbf{Ποιότητα ηλεκτρικής ενέργειας.} Η συχνότητα και οι τάσεις οφείλουν να διατηρούνται εντός
      προδιαγεγραμμένων περιθωρίων, τα οποία για την Ηπειρωτική Ευρώπη ορίζονται
      ρυθμιστικά και όχι κατά την κρίση του Διαχειριστή.
\end{itemize}

Η δυσκολία έγκειται στο ότι οι απαιτήσεις αυτές εκδηλώνονται σε \emph{διαφορετικές χρονικές
κλίμακες}, οι οποίες εκτείνονται σε πάνω από δώδεκα τάξεις μεγέθους. Ο
Πίνακας~\ref{tab:chronikes-klimakes} συνοψίζει την εικόνα και λειτουργεί ως χάρτης για το
υπόλοιπο κεφάλαιο.

\begin{table}[!htbp]
\centering
\begin{tabular}{lll}
\hline
Χρονική κλίμακα & Φαινόμενο ή διαδικασία & Μηχανισμός \\
\hline
μs -- ms   & οδεύοντα κύματα, βραχυκυκλώματα & προστασία, διακόπτες ισχύος \\
ms -- s    & ηλεκτρομηχανικές ταλαντώσεις & αδράνεια, σταθεροποιητές (PSS), FACTS \\
s -- 30 s  & απόκλιση συχνότητας         & εφεδρεία διατήρησης συχνότητας (ΕΔΣ) \\
30 s -- 15 min & αποκατάσταση συχνότητας & εφεδρεία αποκατάστασης συχνότητας (ΕΑΣ) \\
min -- h   & αναπρογραμματισμός          & ενδοημερήσια αγορά, εξισορρόπηση \\
h -- ημέρα & ένταξη μονάδων, κατανομή φορτίου & αγορά επόμενης ημέρας \\
εβδομάδες -- έτη & επάρκεια ισχύος, συντηρήσεις & προγραμματισμός, ανάπτυξη δικτύου \\
\hline
\end{tabular}
\caption{Χρονικές κλίμακες λειτουργίας του ΣΜΗΕ. Σε κάθε γραμμή αντιστοιχεί διαφορετικός
μηχανισμός ελέγχου· το παρόν κεφάλαιο εστιάζει από τα δευτερόλεπτα και άνω.}
\label{tab:chronikes-klimakes}
\end{table}

\subsection{Καταστάσεις λειτουργίας και κριτήριο N-1}

Ένα σύστημα δεν χαρακτηρίζεται απλώς ως ασφαλές ή μη ασφαλές. Η καθιερωμένη σήμερα εικόνα
οφείλεται στον Dy Liacco, ο οποίος στα τέλη της δεκαετίας του 1960 διατύπωσε την αρχή ότι
ο έλεγχος οφείλει να \emph{προσαρμόζεται} στην εκάστοτε κατάσταση ασφάλειας. Ο Κανονισμός (ΕΕ) 2017/1485 για τη λειτουργία του
συστήματος μεταφοράς (System Operation Guideline, SO GL) κωδικοποιεί την αρχή αυτή στο
άρθρο 18, όπου διακρίνει πέντε καταστάσεις του συστήματος:

\begin{enumerate}
\item \textbf{Κανονική κατάσταση} (normal state). Τάσεις, ροές ισχύος και συχνότητα
      βρίσκονται εντός ορίων \emph{και} το σύστημα παραμένει εντός ορίων μετά από
      οποιοδήποτε απρόβλεπτο συμβάν του καταλόγου, λαμβάνοντας υπόψη τα διαθέσιμα
      διορθωτικά μέτρα. Ο έλεγχος είναι \emph{προληπτικός}: ο Διαχειριστής
      προετοιμάζεται για διαταραχές που δεν έχουν εκδηλωθεί.
\item \textbf{Κατάσταση συναγερμού} (alert state· συχνά αποδίδεται και ως κατάσταση
      επιφυλακής). Τα μεγέθη παραμένουν εντός ορίων, όμως τουλάχιστον ένα απρόβλεπτο
      συμβάν του καταλόγου θα οδηγούσε σε παραβίαση ορίων ακόμη και μετά τα διορθωτικά
      μέτρα, ή η εφεδρική δυναμικότητα έχει μειωθεί σημαντικά. Το σύστημα λειτουργεί, αλλά
      μια επόμενη διαταραχή θα το οδηγήσει εκτός ορίων.
\item \textbf{Κατάσταση έκτακτης ανάγκης} (emergency state). Παραβιάζεται τουλάχιστον
      ένα όριο επιχειρησιακής ασφάλειας ή εφαρμόζεται μέτρο του σχεδίου άμυνας του
      συστήματος. Ο έλεγχος καθίσταται \emph{διορθωτικός} και επείγων.
\item \textbf{Κατάσταση γενικής διακοπής} (blackout state· συχνά και «συσκότιση»).
      Απώλεια άνω του 50\% της ζήτησης ή πλήρης απουσία τάσης για τουλάχιστον τρία λεπτά
      στην περιοχή ελέγχου του ΔΣΜ.
\item \textbf{Κατάσταση αποκατάστασης} (restoration state). Ο ΔΣΜ έχει αρχίσει να
      εφαρμόζει τα μέτρα του σχεδίου αποκατάστασης για την επανατροφοδότηση του συστήματος.
\end{enumerate}

Η μετάβαση από την κανονική κατάσταση στην κατάσταση συναγερμού \emph{δεν} συνοδεύεται από ορατή
μεταβολή των μετρούμενων μεγεθών, καθώς κανένα όριο δεν υπερβαίνεται. Πρόκειται για
μεταβολή που υφίσταται μόνο στους υπολογισμούς του Διαχειριστή. Ακριβώς γι' αυτό απαιτείται
αξιόπιστη εικόνα του συστήματος: χωρίς αυτήν, η διαφορά μεταξύ ασφαλούς και επικίνδυνης
λειτουργίας παραμένει αόρατη.

\subsubsection*{Το κριτήριο N-1}

Το \textbf{κριτήριο N-1} αποτελεί τον πρακτικό ορισμό της ασφάλειας: από τα $N$ στοιχεία
του συστήματος, η απώλεια οποιουδήποτε \emph{ενός} — γραμμής, μετασχηματιστή, μονάδας
ηλεκτροπαραγωγής ή ζυγού — πρέπει να μπορεί να απορροφηθεί χωρίς παραβίαση ορίων και χωρίς
αλυσιδωτές αποζεύξεις.

Ο έλεγχός του πραγματοποιείται με \emph{ανάλυση απρόβλεπτων συμβάντων} (contingency
analysis· στη βιβλιογραφία και «ανάλυση ενδεχομένων»): για κάθε συμβάν του καταλόγου
απρόβλεπτων συμβάντων εκτελείται υπολογισμός ροής φορτίου (load flow) με το αντίστοιχο
στοιχείο εκτός λειτουργίας και ελέγχεται τυχόν υπέρβαση ορίων. Σε σύστημα με μερικές
εκατοντάδες κλάδους, αυτό συνεπάγεται εκατοντάδες υπολογισμούς ροής φορτίου, οι οποίοι
πρέπει να ολοκληρωθούν μέσα σε λίγα λεπτά.

Δύο παρατηρήσεις αξίζει να συγκρατηθούν:

\begin{itemize}
\item Το N-1 αποτελεί \emph{ντετερμινιστικό} κριτήριο, καθώς δεν σταθμίζει την πιθανότητα
      κάθε ενδεχομένου. Ένα σπάνιο και ένα συχνό σφάλμα αντιμετωπίζονται ισοδύναμα. Είναι
      απλό και ελέγξιμο, αλλά συντηρητικό σε ορισμένες περιπτώσεις και ανεπαρκές σε άλλες —
      εξ ου και το ενδιαφέρον για πιθανοτικά κριτήρια ασφάλειας.
\item Η ανάλυση απρόβλεπτων συμβάντων απαιτεί ως \emph{είσοδο} την τρέχουσα κατάσταση του συστήματος.
      Αν η είσοδος αυτή είναι εσφαλμένη, εσφαλμένα είναι και όλα τα συμπεράσματα περί
      ασφάλειας. Το γεγονός αυτό καθιστά την εκτίμηση κατάστασης, η οποία εξετάζεται
      παρακάτω, όχι απλώς χρήσιμη αλλά προαπαιτούμενη.
\end{itemize}

\subsection{Εποπτεία και έλεγχος: SCADA και EMS}

\subsubsection*{Η αλυσίδα μέτρησης}

Η βασική υποδομή τηλεμέτρησης και τηλεχειρισμού του ΣΜΗΕ είναι το σύστημα εποπτικού
ελέγχου και συλλογής δεδομένων (\textbf{SCADA}, Supervisory Control and Data Acquisition). Η αλυσίδα από το φυσικό μέγεθος έως την οθόνη
του χειριστή περιλαμβάνει τέσσερις κρίκους:

\begin{enumerate}
\item \textbf{Μετασχηματιστές μέτρησης.} Οι μετασχηματιστές τάσης (VT) και έντασης (CT)
      υποβιβάζουν τα μεγέθη του δικτύου σε τιμές κατάλληλες για τις διατάξεις μέτρησης
      και προστασίας, τυπικά 100 V και 1 ή 5 A στο δευτερεύον τύλιγμα. Εισάγουν το πρώτο
      σφάλμα της αλυσίδας, το οποίο καθορίζεται από την κλάση ακρίβειας, τυπικά 0,2 έως 1
      (σφάλμα λόγου 0,2 έως 1\%).
\item \textbf{Μετατροπείς μέτρησης (transducers) και ευφυείς ηλεκτρονικές συσκευές (IED).}
      Υπολογίζουν τα παράγωγα μεγέθη — ενεργό και άεργο ισχύ, ενεργό τιμή (rms) τάσης και
      έντασης — από τα στιγμιαία δείγματα.
\item \textbf{Απομακρυσμένες τερματικές μονάδες (RTU).} Συγκεντρώνουν τα δεδομένα του
      υποσταθμού και τα αποστέλλουν στο Κέντρο Ελέγχου Ενέργειας.
\item \textbf{Τηλεπικοινωνιακό δίκτυο.} Μεταφέρει τα δεδομένα με τυποποιημένα πρωτόκολλα,
      με πλέον διαδεδομένα τα IEC 60870-5-101/104 και το DNP3, ενώ εντός του υποσταθμού
      επικρατεί σήμερα το IEC 61850.
\end{enumerate}

Συλλέγονται δύο κατηγορίες δεδομένων:

\begin{itemize}
\item \textbf{Αναλογικές μετρήσεις} (τηλεμετρήσεις): ροές ενεργού και αέργου ισχύος στα
      άκρα των γραμμών και των μετασχηματιστών, εγχύσεις ισχύος στους ζυγούς, μέτρα τάσεων
      ζυγών, εντάσεις ρεύματος κλάδων.
\item \textbf{Ψηφιακές ενδείξεις} (τηλεσημάνσεις): θέσεις διακοπτών ισχύος και αποζευκτών,
      θέσεις των μεταγωγέων λήψεων υπό φορτίο των μετασχηματιστών, σήματα συναγερμού και
      λειτουργίας των ηλεκτρονόμων προστασίας. Οι ενδείξεις αυτές
      καθορίζουν την \emph{τοπολογία}, δηλαδή ποια στοιχεία βρίσκονται εν λειτουργία και
      με ποιον τρόπο συνδέονται μεταξύ τους.
\end{itemize}

Η περίοδος σάρωσης ενός συμβατικού συστήματος SCADA είναι της τάξης λίγων δευτερολέπτων,
συνήθως 2 έως 10 s.

\subsubsection*{Δύο καθοριστικοί περιορισμοί}

Δύο χαρακτηριστικά του SCADA είναι κρίσιμα για όσα ακολουθούν και αξίζει να συγκρατηθούν:

\begin{itemize}
\item Οι μετρήσεις \textbf{δεν είναι χρονικά συγχρονισμένες} μεταξύ τους. Συλλέγονται σε
      διαφορετικές χρονικές στιγμές εντός του κύκλου σάρωσης, οπότε το σύνολό τους αποτελεί
      \emph{οιονεί στατική} απεικόνιση και όχι στιγμιότυπο. Ένα σύνολο μετρήσεων που έχει
      κατανεμηθεί σε πέντε δευτερόλεπτα δεν αντιστοιχεί σε καμία πραγματική στιγμή του
      συστήματος. Το γεγονός αυτό είναι αμελητέο όσο το σύστημα μεταβάλλεται αργά, αλλά
      καθίσταται σοβαρό κατά τη διάρκεια διαταραχής.
\item Οι μετρήσεις είναι \textbf{θορυβώδεις και ενίοτε εσφαλμένες}, ενώ ορισμένα μεγέθη δεν
      μετρώνται καθόλου. Χαρακτηριστικά, οι \emph{γωνίες} των τάσεων — οι οποίες αποτελούν
      το ήμισυ της κατάστασης του συστήματος — απουσιάζουν εξ ολοκλήρου από το συμβατικό
      SCADA.
\end{itemize}

Ένας μετασχηματιστής έντασης με εσφαλμένα καταχωρημένο λόγο μετασχηματισμού, μια μέτρηση
ισχύος με ανεστραμμένη πολικότητα, ένας διακόπτης του οποίου η ένδειξη θέσης δεν
ανανεώνεται λόγω βλάβης: όλες αυτές οι περιπτώσεις
παράγουν δεδομένα που εμφανίζονται απολύτως φυσιολογικά. Το σύστημα εποπτείας δεν διαθέτει
τρόπο να τα διακρίνει \emph{ανά μέτρηση}. Η διάκριση είναι εφικτή μόνο συλλογικά, με
αξιοποίηση του πλεονασμού των μετρήσεων (redundancy), όπως αναλύεται στη συνέχεια.

\subsubsection*{Το EMS και η σειρά των λειτουργιών}

Το σύστημα διαχείρισης ενέργειας (\textbf{EMS}, Energy Management System) είναι το λογισμικό που μετατρέπει τα ακατέργαστα
δεδομένα σε αποφάσεις. Η σειρά εκτέλεσης των λειτουργιών του δεν είναι αυθαίρετη, καθώς
κάθε βήμα τροφοδοτεί το επόμενο:

\begin{enumerate}
\item \textbf{Επεξεργαστής τοπολογίας} (topology processor). Συνθέτει, από τις θέσεις των
      διακοπτών και των αποζευκτών, το μοντέλο ζυγών-κλάδων (bus-branch model) του δικτύου. Προσδιορίζει τον
      τρόπο με τον οποίο είναι συνδεδεμένο το δίκτυο τη δεδομένη στιγμή.
\item \textbf{Εκτίμηση κατάστασης}. Προσδιορίζει την ηλεκτρική κατάσταση του δικτύου αυτού.
\item \textbf{Ανάλυση απρόβλεπτων συμβάντων}. Αποφαίνεται για την επιχειρησιακή ασφάλεια
      της κατάστασης αυτής.
\item \textbf{Βέλτιστη ροή φορτίου} (Optimal Power Flow, OPF). Προσδιορίζει, εφόσον
      απαιτείται, τα αναγκαία διορθωτικά μέτρα (remedial actions) και το κόστος τους.
\item \textbf{Αυτόματη Ρύθμιση Παραγωγής} (ΑΡΠ· Automatic Generation Control, AGC).
      Υλοποιεί συνεχώς τη ρύθμιση φορτίου-συχνότητας και την τήρηση των προγραμματισμένων
      ανταλλαγών ισχύος με τις γειτονικές περιοχές.
\end{enumerate}

Σφάλμα στο πρώτο ή στο δεύτερο βήμα διαδίδεται σιωπηρά σε όλα τα επόμενα.

Το Σχήμα~\ref{fig:scada-ems} συνοψίζει την πλήρη διαδρομή. Αξίζει να παρατηρηθεί ότι
πρόκειται για \emph{κλειστό} βρόχο: τα δεδομένα ανεβαίνουν από το δίκτυο προς το κέντρο
ελέγχου και οι εντολές επιστρέφουν σε αυτό.

\begin{figure}[!htbp]
\centering
\includegraphics[width=\linewidth]{figures/scada-ems.svg}
\caption{Ο βρόχος εποπτείας και ελέγχου. Η επάνω σειρά είναι η αλυσίδα μέτρησης, από το
φυσικό δίκτυο έως το SCADA· η κάτω σειρά είναι η ακολουθία λειτουργιών του EMS, η οποία
καταλήγει σε εντολές προς το δίκτυο.}
\label{fig:scada-ems}
\end{figure}

\subsection{Συγχρονισμένες μετρήσεις φασιθετών και συστήματα WAMS}

\subsubsection*{Η σημασία της γωνίας}

Υπενθυμίζεται η σχέση μεταφοράς ενεργού ισχύος μεταξύ δύο ζυγών που συνδέονται μέσω
γραμμής χωρίς απώλειες με επαγωγική αντίδραση σειράς $X$:

\begin{equation}
P_{12} = \frac{V_1 V_2}{X}\,\sin(\delta_1 - \delta_2)
\label{eq:powerflow}
\end{equation}

Η ροή ισχύος καθορίζεται από τη \emph{διαφορά γωνιών}. Η γωνία δεν αποτελεί επομένως
βοηθητικό μαθηματικό μέγεθος, αλλά την κινητήρια αιτία της ροής ισχύος και τον πλέον άμεσο
δείκτη φόρτισης του δικτύου. Η αδυναμία μέτρησής της από το συμβατικό SCADA συνιστά
σοβαρή έλλειψη.

\subsubsection*{Η αρχή λειτουργίας}

Οι \textbf{μονάδες μέτρησης φασιθετών} (Phasor Measurement Units, PMU) αίρουν τον
περιορισμό αυτόν. Μια PMU υπολογίζει τον φασιθέτη — μέτρο \emph{και} γωνία — της τάσης και
της έντασης του ρεύματος και τον συνοδεύει με χρονοσήμανση ως προς την κοινή χρονική αναφορά
UTC, η οποία λαμβάνεται τυπικά μέσω GPS.

Το ουσιώδες σημείο έγκειται στο εξής: η γωνία ενός φασιθέτη έχει νόημα μόνο ως προς κάποια
αναφορά. Επειδή όλες οι PMU ανάγονται στην ίδια χρονική βάση, οι γωνίες που μετρούν είναι
άμεσα συγκρίσιμες σε ολόκληρο το σύστημα, ακόμη και σε αποστάσεις εκατοντάδων χιλιομέτρων —
εξ ου και ο όρος \emph{συγχρονισμένοι φασιθέτες} (synchrophasors). Για να είναι έγκυρη η σύγκριση
σε σύστημα 50 Hz απαιτείται ακρίβεια χρονισμού καλύτερη του μικροδευτερολέπτου: σφάλμα 1 μs
αντιστοιχεί σε σφάλμα γωνίας $360^\circ \times 50 \times 10^{-6} = 0{,}018^\circ$.

\begin{figure}[!htbp]
\centering
\includegraphics[width=\linewidth]{figures/pmu-wams.svg}
\caption{Αριστερά, το μετρούμενο μέγεθος: ο φασιθέτης τάσης με μέτρο και γωνία ως προς
κοινή χρονική αναφορά. Δεξιά, η ιεραρχία συλλογής: οι PMU συγχρονίζονται μέσω GPS, τα
πλαίσια δεδομένων ευθυγραμμίζονται χρονικά στον συγκεντρωτή δεδομένων (PDC) και
τροφοδοτούν το WAMS.}
\label{fig:pmu-wams}
\end{figure}

\subsubsection*{Προδιαγραφές}

Τα βασικά τεχνικά χαρακτηριστικά καθορίζονται από το πρότυπο IEC/IEEE 60255-118-1:2018, το
οποίο διαδέχθηκε τη σειρά IEEE C37.118:

\begin{itemize}
\item \textbf{Ρυθμός αναφοράς} (reporting rate): τυπικά 10, 25 ή 50 πλαίσια δεδομένων ανά
      δευτερόλεπτο για σύστημα
      50 Hz, με δυνατότητα υψηλότερων ρυθμών.
\item \textbf{Ακρίβεια}: το ολικό διανυσματικό σφάλμα (Total Vector Error, TVE) δεν
      επιτρέπεται να υπερβαίνει το 1\% στη μόνιμη κατάσταση. Ορίζεται ως το μέτρο της
      διαφοράς του μετρούμενου από τον πραγματικό φασιθέτη, κανονικοποιημένο ως προς τον
      πραγματικό, και συνεπώς ενοποιεί σε ένα μέγεθος τα σφάλματα μέτρου και γωνίας.
      Ορίζονται επιπλέον το σφάλμα συχνότητας (FE) και το σφάλμα ρυθμού μεταβολής
      συχνότητας (RFE).
\item \textbf{Κλάσεις}: η κλάση P (protection) βελτιστοποιείται για ταχεία απόκριση, με
      μικρή καθυστέρηση φίλτρου, ενώ η κλάση M (measurement) για ακρίβεια υπό την παρουσία
      αρμονικών και παρεμβολών. Η επιλογή συνιστά συμβιβασμό, καθώς δεν υφίσταται φίλτρο
      ταυτόχρονα ταχύ και με ισχυρή απόρριψη παρεμβολών.
\end{itemize}

\begin{table}[!htbp]
\centering
\begin{tabular}{lll}
\hline
 & SCADA & PMU / WAMS \\
\hline
Ανανέωση         & ανά 2--10 s        & 10--50 πλαίσια ανά s \\
Συγχρονισμός     & όχι                & ναι (UTC μέσω GPS) \\
Γωνία τάσης      & δεν μετράται       & μετράται άμεσα \\
Σχέση με την κατάσταση & μη γραμμική & γραμμική \\
Κάλυψη           & εκτεταμένη         & επιλεκτική, λόγω κόστους \\
Τυπική χρήση     & εκτίμηση κατάστασης, τηλεχειρισμοί & εποπτεία δυναμικών φαινομένων \\
\hline
\end{tabular}
\caption{Σύγκριση συμβατικής και συγχρονισμένης μέτρησης. Οι δύο τεχνολογίες είναι
συμπληρωματικές και όχι ανταγωνιστικές.}
\label{tab:scada-pmu}
\end{table}

\subsubsection*{Από τη μέτρηση στην εφαρμογή}

Οι μετρήσεις συγκεντρώνονται ιεραρχικά σε συγκεντρωτές δεδομένων φασιθετών (Phasor Data
Concentrators, PDC), οι οποίοι ευθυγραμμίζουν χρονικά τα πλαίσια από διαφορετικές PMU και
συνθέτουν ένα \textbf{σύστημα εποπτείας ευρείας περιοχής} (Wide Area Monitoring System,
WAMS). Οι χαρακτηριστικές εφαρμογές είναι:

\begin{itemize}
\item η παρακολούθηση ηλεκτρομηχανικών ταλαντώσεων μεταξύ περιοχών (inter-area
      oscillations), με συχνότητες τυπικά
      0,1 έως 1 Hz, οι οποίες παραμένουν πρακτικά αόρατες σε ρυθμό σάρωσης δευτερολέπτων·
\item η μέτρηση του ρυθμού μεταβολής της συχνότητας και η εκτίμηση της διαθέσιμης
      αδράνειας·
\item η ανίχνευση νησιδοποίησης, δηλαδή της μετάβασης τμήματος του δικτύου σε απομονωμένη
      λειτουργία, και η επιτήρηση της γωνιακής απόκλισης μεταξύ περιοχών·
\item η επικύρωση μοντέλων μετά από διαταραχές, μέσω σύγκρισης της καταγεγραμμένης με την
      προσομοιωμένη απόκριση·
\item η ενίσχυση — ή, σε συνθήκες πλήρους κάλυψης, η αντικατάσταση — της κλασικής εκτίμησης
      κατάστασης.
\end{itemize}

\subsection{Ευέλικτα συστήματα μεταφοράς εναλλασσόμενου ρεύματος (FACTS)}

Ενώ τα συστήματα SCADA και PMU \emph{παρατηρούν} το σύστημα, οι διατάξεις \textbf{FACTS}
(Flexible AC Transmission Systems) το \emph{ελέγχουν} ενεργά.

Η \eqref{eq:powerflow} προσδιορίζει με ακρίβεια τα ελεγχόμενα μεγέθη. Η ροή ισχύος
εξαρτάται από τρία και μόνο μεγέθη: τις τάσεις $V_1, V_2$, την αντίδραση $X$ και τη διαφορά
γωνιών $\delta_1 - \delta_2$. Σε κλασικό δίκτυο εναλλασσόμενου ρεύματος και τα τρία είναι
ουσιαστικά δεδομένα, οπότε η ισχύς ρέει σύμφωνα με τους νόμους της φυσικής και όχι σύμφωνα
με τις επιδιώξεις του Διαχειριστή. Οι διατάξεις FACTS είναι μετατροπείς ηλεκτρονικών ισχύος
που καθιστούν τα μεγέθη αυτά ελεγχόμενα, ταχέως και συνεχώς. Κατατάσσονται ανάλογα με το
μέγεθος που επηρεάζουν, όπως συνοψίζει το Σχήμα~\ref{fig:facts}:

\begin{itemize}
\item \textbf{Παράλληλης σύνδεσης} (έλεγχος του $V$): ο στατός αντισταθμιστής αέργου ισχύος
      (Static Var Compensator, SVC) και ο στατός σύγχρονος αντισταθμιστής (Static
      Synchronous Compensator, STATCOM). Ελέγχουν την τάση του ζυγού μέσω έγχυσης
      ή απορρόφησης αέργου ισχύος.
\item \textbf{Σειριακής σύνδεσης} (έλεγχος του $X$): ο πυκνωτής σειράς ελεγχόμενος με
      θυρίστορ (Thyristor Controlled Series Capacitor, TCSC) και ο στατός σύγχρονος
      αντισταθμιστής σειράς (Static Synchronous Series Compensator, SSSC). Μεταβάλλοντας
      την ισοδύναμη αντίδραση σειράς της γραμμής, ανακατανέμουν τις ροές ισχύος και
      αυξάνουν τη μεταφορική ικανότητα.
\item \textbf{Συνδυασμένης σύνδεσης}, σε σειρά και παράλληλα (έλεγχος και της γωνίας): ο
      ενοποιημένος ελεγκτής ροής ισχύος (Unified Power Flow Controller, UPFC) ελέγχει
      ταυτόχρονα και ανεξάρτητα την τάση και τις ροές ενεργού και αέργου ισχύος. Αποτελεί την πληρέστερη και ταυτόχρονα την ακριβότερη
      λύση.
\end{itemize}

\begin{figure}[!htbp]
\centering
\includegraphics[width=0.88\linewidth]{figures/facts.svg}
\caption{Οι τρεις κατηγορίες διατάξεων FACTS σε γραμμή δύο ζυγών. Κάθε κατηγορία επεμβαίνει
σε διαφορετικό μέγεθος της σχέσης μεταφοράς ισχύος: η παράλληλη στην τάση, η σειριακή στην
αντίδραση, η συνδυασμένη και στη γωνία.}
\label{fig:facts}
\end{figure}

\subsubsection*{Εσωτερική δομή}

Η ταξινόμηση κατά τρόπο σύνδεσης δεν εξαντλεί το θέμα. Καθοριστική είναι και μια δεύτερη
διάκριση, εκείνη ανάμεσα σε διατάξεις που συμπεριφέρονται ως \emph{μεταβλητή σύνθετη αντίσταση} (εμπέδηση) και
σε διατάξεις που συμπεριφέρονται ως \emph{ελεγχόμενη πηγή τάσης}. Το
Σχήμα~\ref{fig:facts-devices} δίνει τα μονογραμμικά διαγράμματα.

\begin{itemize}
\item Ο \textbf{SVC} αποτελείται από πηνίο ελεγχόμενο με θυρίστορ (Thyristor Controlled
      Reactor, TCR) και πυκνωτές μεταγόμενους με θυρίστορ (Thyristor Switched Capacitor,
      TSC), σε παράλληλη σύνδεση. Η γωνία έναυσης του TCR ρυθμίζει συνεχώς την ισοδύναμη
      επαγωγική επιδεκτικότητα, ενώ οι TSC προσθέτουν χωρητική επιδεκτικότητα σε διακριτά
      βήματα. Το σύνολο συμπεριφέρεται ως \emph{μεταβλητή επιδεκτικότητα} $B$ (susceptance).
\item Ο \textbf{STATCOM} είναι μετατροπέας πηγής τάσης (Voltage Source Converter, VSC) με
      πυκνωτή στην πλευρά συνεχούς ρεύματος (ΣΡ). Παράγει εναλλασσόμενη τάση ελεγχόμενου πλάτους·
      η διαφορά της από την τάση του δικτύου, εφαρμοζόμενη στην αντίδραση σύζευξης,
      καθορίζει το μέτρο και τη φορά της αέργου ισχύος.
\item Ο \textbf{TCSC} είναι πυκνωτής σειράς με παράλληλο TCR, οπότε η ισοδύναμη αντίδραση
      σειράς της γραμμής μεταβάλλεται συνεχώς. Ο \textbf{SSSC} εγχέει αντ' αυτού τάση σε
      σειρά μέσω μετασχηματιστή σύζευξης σειράς: η εγχεόμενη τάση είναι ανεξάρτητη του ρεύματος γραμμής, ενώ
      στον TCSC είναι κατ' ανάγκη ανάλογη αυτού.
\item Ο \textbf{UPFC} συνδυάζει έναν παράλληλο και έναν σειριακό VSC σε \emph{κοινό ζυγό ΣΡ}
      (DC link). Ο κοινός αυτός ζυγός είναι το κρίσιμο στοιχείο: επιτρέπει την ανταλλαγή
      \emph{ενεργού} ισχύος μεταξύ των δύο μετατροπέων, ώστε ο σειριακός να εγχέει τάση
      οποιουδήποτε πλάτους και οποιασδήποτε φάσης — και συνεπώς να ελέγχει ανεξάρτητα
      τις ροές ενεργού και αέργου ισχύος.
\end{itemize}

\begin{figure}[!htbp]
\centering
\includegraphics[width=\linewidth]{figures/facts-devices.svg}
\caption{Μονογραμμικά διαγράμματα των βασικών διατάξεων FACTS. Οι δύο πρώτες συνδέονται
παράλληλα, οι δύο επόμενες σε σειρά, ενώ ο UPFC συνδυάζει αμφότερους τους τρόπους μέσω
κοινού ζυγού ΣΡ.}
\label{fig:facts-devices}
\end{figure}

\subsubsection*{Γιατί υπερτερεί ο STATCOM στη χαμηλή τάση}

Η διάκριση μεταξύ μεταβλητής εμπέδησης και πηγής τάσης δεν είναι ακαδημαϊκή: αποτυπώνεται
άμεσα στα χαρακτηριστικά $V$--$I$ των δύο διατάξεων, τα οποία συγκρίνει το
Σχήμα~\ref{fig:facts-vi}.

Ο SVC, ως μεταβλητή επιδεκτικότητα, υπακούει στη σχέση $I = B\,V$. Ακόμη και με πλήρως ενεργοποιημένη
τη χωρητική του ικανότητα, το ρεύμα που αποδίδει είναι \emph{ανάλογο} της τάσης του ζυγού:
στο μισό της ονομαστικής τάσης αποδίδει το μισό ρεύμα, και επομένως το ένα τέταρτο της
αέργου ισχύος. Ο STATCOM αντιθέτως, ως πηγή τάσης με έλεγχο ρεύματος, διατηρεί το πλήρες
ονομαστικό ρεύμα έως πολύ χαμηλές τιμές τάσης, περιοριζόμενος μόνο από το μέγιστο επιτρεπόμενο
ρεύμα των ημιαγωγικών διακοπτών του.

Η διαφορά αυτή εκδηλώνεται ακριβώς όταν έχει τη μεγαλύτερη σημασία, δηλαδή αμέσως μετά από
σφάλμα, όταν η τάση έχει καταρρεύσει και απαιτείται μέγιστη στήριξη. Σε αυτό οφείλεται η
υπεροχή του STATCOM στη δυναμική στήριξη τάσης.

\begin{figure}[!htbp]
\centering
\includegraphics[width=\linewidth]{figures/facts-vi.svg}
\caption{Χαρακτηριστικά $V$--$I$ των δύο διατάξεων παράλληλης σύνδεσης. Τα όρια του SVC
είναι ευθείες σταθερής επιδεκτικότητας που διέρχονται από την αρχή, ενώ του STATCOM είναι
κατακόρυφες ευθείες σταθερού ρεύματος.}
\label{fig:facts-vi}
\end{figure}

Σε συνδυασμό με τα συστήματα συνεχούς ρεύματος υψηλής τάσης (HVDC), όπου η ροή ισχύος
ορίζεται απευθείας από τον έλεγχο του μετατροπέα, οι διατάξεις FACTS αποτελούν τα βασικά
εργαλεία του Διαχειριστή για τη ρύθμιση τάσης και τη διαχείριση συμφορήσεων στο σύστημα
μεταφοράς. Τα αντίστοιχα εργαλεία στα δίκτυα διανομής εξετάζονται στα Κεφάλαια 5 και 6.

\subsection{Εκτίμηση κατάστασης}

\subsubsection*{Διατύπωση του προβλήματος}

Όπως αναλύθηκε παραπάνω, οι μετρήσεις του SCADA είναι πλεονάζουσες, θορυβώδεις, ασύγχρονες
και ελλιπείς, ενώ παράλληλα η ανάλυση ασφάλειας απαιτεί \emph{ενιαία και συνεπή} εικόνα του
δικτύου. Η \textbf{εκτίμηση κατάστασης} (state estimation) γεφυρώνει το χάσμα αυτό.
Εισήχθη στα συστήματα ηλεκτρικής ενέργειας από τους Schweppe και Wildes το 1970, σε σειρά
τριών εργασιών οι οποίες παραμένουν έως σήμερα σημείο αναφοράς, και αποτελεί έκτοτε τη
θεμελιώδη λειτουργία κάθε EMS.

Μια πρώτη προσέγγιση θα ήταν ένας απλός υπολογισμός ροής φορτίου. Η προσέγγιση αυτή δεν
επαρκεί, για δύο λόγους. Πρώτον, η ροή φορτίου απαιτεί \emph{ακριβώς} τα κατάλληλα δεδομένα
εισόδου, ενώ τα διαθέσιμα δεδομένα είναι πολύ περισσότερα από τα αναγκαία και όλα ελαφρώς
εσφαλμένα. Δεύτερον, η ροή φορτίου δεν διαθέτει μηχανισμό ανίχνευσης εσφαλμένων δεδομένων
και θα συγκλίνει απρόσκοπτα σε εσφαλμένη λύση. Απαιτείται επομένως μέθοδος που
\emph{αξιοποιεί} την πλεονάζουσα πληροφορία αντί να την αγνοεί.

Το Σχήμα~\ref{fig:se-flow} δίνει τη συνολική ροή του αλγορίθμου. Οι επιμέρους βαθμίδες
αναλύονται στις υποενότητες που ακολουθούν· χρήσιμο είναι να σημειωθεί από τώρα ότι
υπάρχουν \emph{δύο} εμφωλευμένοι βρόχοι — ο εσωτερικός για τη σύγκλιση της αριθμητικής
επίλυσης και ο εξωτερικός για την απόρριψη εσφαλμένων μετρήσεων.

\begin{figure}[!htbp]
\centering
\includegraphics[width=0.8\linewidth]{figures/se-flowchart.svg}
\caption{Ροή της εκτίμησης κατάστασης σε ένα EMS. Ο εσωτερικός βρόχος (αριστερά) επαναλαμβάνει
τη μέθοδο Gauss--Newton έως τη σύγκλιση· ο εξωτερικός (δεξιά) απορρίπτει τη μέτρηση με το
μέγιστο κανονικοποιημένο υπόλοιπο και επανεκτιμά.}
\label{fig:se-flow}
\end{figure}

\subsubsection*{Το μοντέλο μέτρησης}

Ως \emph{κατάσταση} ορίζεται το διάνυσμα
$\mathbf{x} = [\boldsymbol{\theta}^{T}, \mathbf{V}^{T}]^{T}$ των γωνιών και των μέτρων των
τάσεων όλων των ζυγών. Για δίκτυο $N$ ζυγών, η γωνία ενός ζυγού, του ζυγού αναφοράς
(reference bus), τίθεται αυθαίρετα, συνήθως ίση με μηδέν, οπότε το πλήθος των αγνώστων ανέρχεται σε

\begin{equation}
n = 2N - 1
\label{eq:nstates}
\end{equation}

Η επιλογή αυτή είναι θεμελιώδης: εφόσον είναι γνωστά τα $\mathbf{V}$ και
$\boldsymbol{\theta}$, \emph{κάθε} άλλο μέγεθος του δικτύου — ροές, εγχύσεις, απώλειες —
υπολογίζεται από τις εξισώσεις ροής φορτίου. Η κατάσταση αποτελεί το ελάχιστο σύνολο μεγεθών που
περιγράφει πλήρως το σύστημα.

Το μοντέλο μέτρησης δίνεται από τη σχέση

\begin{equation}
\mathbf{z} = \mathbf{h}(\mathbf{x}) + \mathbf{e}
\label{eq:measmodel}
\end{equation}

όπου $\mathbf{z} \in \mathbb{R}^{m}$ το διάνυσμα των μετρήσεων, $\mathbf{h}(\cdot)$ οι μη
γραμμικές συναρτήσεις που συνδέουν την κατάσταση με τα μετρούμενα μεγέθη, και $\mathbf{e}$
ο θόρυβος μέτρησης. Ο τελευταίος θεωρείται μηδενικής μέσης τιμής, με διαγώνιο πίνακα
συνδιακύμανσης $\mathbf{R} = \mathrm{diag}(\sigma_1^2, \ldots, \sigma_m^2)$, δηλαδή τα
σφάλματα των επιμέρους μετρητών θεωρούνται ασυσχέτιστα.

Καθοριστικό μέγεθος αποτελεί ο \textbf{λόγος πλεονασμού}

\begin{equation}
\eta = \frac{m}{n}
\label{eq:redundancy}
\end{equation}

ο οποίος σε πραγματικά συστήματα κυμαίνεται τυπικά μεταξύ 1,7 και 3. Ενδεικτικά, σε δίκτυο
$N = 100$ ζυγών οι άγνωστοι ανέρχονται σε $n = 199$ και, για $\eta = 2$, οι μετρήσεις σε
$m = 400$ περίπου. Ο πλεονασμός είναι εκείνος που επιτρέπει τόσο τη μείωση του θορύβου όσο
και, το κυριότερο, την ανίχνευση εσφαλμένων μετρήσεων. Ελλείψει αυτού, το πρόβλημα
εκφυλίζεται σε υπολογισμό ροής φορτίου, με όλα τα μειονεκτήματα που προαναφέρθηκαν.

\subsubsection*{Σταθμισμένα ελάχιστα τετράγωνα}

Επειδή οι μετρήσεις είναι περισσότερες από τους αγνώστους, δεν υφίσταται $\mathbf{x}$ που
να τις ικανοποιεί όλες. Αναζητείται επομένως εκείνο που τις ικανοποιεί κατά τον βέλτιστο
δυνατό τρόπο. Η καθιερωμένη διατύπωση είναι εκείνη των \textbf{σταθμισμένων ελαχίστων
τετραγώνων} (Weighted Least Squares, WLS):

\begin{equation}
J(\mathbf{x}) = \sum_{i=1}^{m} \frac{\left(z_i - h_i(\mathbf{x})\right)^2}{\sigma_i^2}
             = \left[\mathbf{z}-\mathbf{h}(\mathbf{x})\right]^{T} \mathbf{R}^{-1}
               \left[\mathbf{z}-\mathbf{h}(\mathbf{x})\right]
\label{eq:wls}
\end{equation}

Η στάθμιση με $1/\sigma_i^2$ έχει άμεση φυσική σημασία: μια μέτρηση με μικρή τυπική
απόκλιση, δηλαδή ακριβής, βαρύνει περισσότερο, ενώ μια θορυβώδης βαρύνει λιγότερο.
Ισοδύναμα, κάθε όρος του αθροίσματος αποτελεί το τετράγωνο του σφάλματος \emph{εκφρασμένο
σε μονάδες της οικείας τυπικής απόκλισης}, ώστε να καθίστανται συγκρίσιμα μεταξύ τους
μεγέθη διαφορετικής φύσης και κλίμακας. Υπό την υπόθεση κανονικά κατανεμημένων σφαλμάτων,
η ελαχιστοποίηση της \eqref{eq:wls} ταυτίζεται με την εκτίμηση μέγιστης πιθανοφάνειας.

Επειδή η $\mathbf{h}$ είναι μη γραμμική, η λύση προκύπτει επαναληπτικά με τη μέθοδο
Gauss--Newton, μέσω διαδοχικών βημάτων γραμμικοποίησης γύρω από την τρέχουσα εκτίμηση,
επίλυσης του γραμμικού προβλήματος και ενημέρωσης. Με
$\mathbf{H}(\mathbf{x}) = \partial \mathbf{h}/\partial \mathbf{x}$ τον ιακωβιανό πίνακα, σε
κάθε επανάληψη $k$ επιλύεται το σύστημα των \emph{κανονικών εξισώσεων}

\begin{align}
\mathbf{G}(\mathbf{x}^{k}) \, \Delta \mathbf{x}^{k}
  &= \mathbf{H}^{T}(\mathbf{x}^{k}) \, \mathbf{R}^{-1}
     \left[\mathbf{z} - \mathbf{h}(\mathbf{x}^{k})\right] \\
\mathbf{G}(\mathbf{x}^{k})
  &= \mathbf{H}^{T}(\mathbf{x}^{k}) \, \mathbf{R}^{-1} \mathbf{H}(\mathbf{x}^{k})
\label{eq:gain}
\end{align}

και ενημερώνεται η εκτίμηση ως $\mathbf{x}^{k+1} = \mathbf{x}^{k} + \Delta \mathbf{x}^{k}$,
έως ότου το $\max|\Delta \mathbf{x}^{k}|$ κατέλθει κάτω από προκαθορισμένη ανοχή
σύγκλισης. Η σύγκλιση είναι τυπικά ταχεία, της τάξης των τριών έως πέντε επαναλήψεων,
επειδή ως αρχική εκτίμηση χρησιμοποιείται συνήθως η λύση της προηγούμενης εκτέλεσης, ενώ το
σύστημα δεν έχει μεταβληθεί σημαντικά στο μεσοδιάστημα.

Ο πίνακας $\mathbf{G}$ ονομάζεται \emph{πίνακας κέρδους} (gain matrix). Είναι συμμετρικός,
θετικά ορισμένος όταν το σύστημα είναι παρατηρήσιμο και, το κυριότερο, \emph{αραιός}, καθώς κάθε μέτρηση εμπλέκει ελάχιστους
μόνο ζυγούς. Η αραιότητα δεν συνιστά λεπτομέρεια υλοποίησης: είναι εκείνη που καθιστά το
πρόβλημα επιλύσιμο σε πραγματικό χρόνο για δίκτυα χιλιάδων ζυγών. Μια πυκνή παραγοντοποίηση
θα απαιτούσε χρόνο ανάλογο του $n^3$, ενώ η αραιή κλιμακώνεται σχεδόν γραμμικά.

\subsubsection*{Παρατηρησιμότητα}

Η επίλυση της \eqref{eq:gain} προϋποθέτει αντιστρεψιμότητα του $\mathbf{G}$, η οποία δεν
είναι δεδομένη: μεγάλο πλήθος μετρήσεων δεν συνεπάγεται κατ' ανάγκη \emph{κατάλληλες}
μετρήσεις. Εφόσον όλες συγκεντρώνονται σε ένα τμήμα του δικτύου, το υπόλοιπο παραμένει
άγνωστο ανεξαρτήτως πλήθους.

Το σύστημα χαρακτηρίζεται \textbf{παρατηρήσιμο} όταν οι διαθέσιμες μετρήσεις επαρκούν για
τον μονοσήμαντο προσδιορισμό ολόκληρης της κατάστασης. Ο έλεγχος πραγματοποιείται με δύο
ισοδύναμες προσεγγίσεις:

\begin{itemize}
\item \textbf{Τοπολογική}: αναζητείται επικαλύπτον δένδρο (spanning tree) του δικτύου πλήρως
      καλυμμένο από μετρήσεις. Είναι ταχεία και προσφέρει άμεση φυσική εποπτεία.
\item \textbf{Αριθμητική}: ελέγχεται η τάξη (rank) του $\mathbf{H}$ ή η παραγοντοποίηση του
      $\mathbf{G}$. Είναι γενικότερη και εντοπίζει και τις οριακές περιπτώσεις.
\end{itemize}

Όταν το σύστημα δεν είναι πλήρως παρατηρήσιμο, εντοπίζονται οι \emph{παρατηρήσιμες νησίδες}
και η κατάσταση εκτιμάται μόνο εντός αυτών. Τα κενά γεφυρώνονται με \emph{ψευδομετρήσεις}
(pseudo-measurements), δηλαδή εκτιμήσεις φορτίου από ιστορικά δεδομένα ή προβλέψεις, στις
οποίες αποδίδεται μεγάλη τυπική απόκλιση ώστε να βαρύνουν ελάχιστα.

\subsubsection*{Ανίχνευση και αναγνώριση εσφαλμένων μετρήσεων}

Στο σημείο αυτό αξιοποιείται ο πλεονασμός. Εσφαλμένη μέτρηση είναι ασύμβατη με τις
υπόλοιπες, και η ασυμβατότητα αυτή αποτυπώνεται στα \emph{υπόλοιπα} (residuals)
$\mathbf{r} = \mathbf{z}-\mathbf{h}(\hat{\mathbf{x}})$. Η κλασική διαδικασία περιλαμβάνει
δύο στάδια.

\textbf{Στάδιο 1 --- Ανίχνευση.} Υπό την υπόθεση ότι όλες οι μετρήσεις είναι ορθές, η
αντικειμενική συνάρτηση $J(\hat{\mathbf{x}})$ στη λύση ακολουθεί κατανομή $\chi^2$ με
$m-n$ βαθμούς ελευθερίας. Στο προηγούμενο αριθμητικό παράδειγμα, με $m = 400$ και
$n = 199$, οι βαθμοί ελευθερίας είναι 201 και η αναμενόμενη τιμή του $J$ επίσης περίπου
201. Υπέρβαση του κατωφλίου $\chi^2_{m-n,\,p}$, για πιθανότητα $p$ συνήθως ίση με 95\%,
συνιστά στατιστική ένδειξη ύπαρξης εσφαλμένων δεδομένων (bad data). Ο έλεγχος ωστόσο
αποφαίνεται μόνο για την \emph{ύπαρξη} του προβλήματος, όχι για τον εντοπισμό του.

\textbf{Στάδιο 2 --- Αναγνώριση.} Τα ακατέργαστα υπόλοιπα δεν είναι συγκρίσιμα μεταξύ τους,
καθώς κάθε μέτρηση χαρακτηρίζεται από διαφορετική ακρίβεια και διαφορετικό βαθμό στήριξης
από τις γειτονικές. Κανονικοποιούνται επομένως ως

\begin{equation}
r_i^{N} = \frac{|r_i|}{\sqrt{\Omega_{ii}}}, \qquad
\boldsymbol{\Omega} = \mathbf{R} - \mathbf{H}\mathbf{G}^{-1}\mathbf{H}^{T}
\label{eq:normres}
\end{equation}

όπου $\boldsymbol{\Omega}$ ο πίνακας συνδιακύμανσης των υπολοίπων. Η μέτρηση με το μέγιστο
$r_i^{N}$, σύμφωνα με το κριτήριο του μεγίστου κανονικοποιημένου υπολοίπου (Largest
Normalized Residual, LNR), θεωρείται ύποπτη, απορρίπτεται, και η εκτίμηση επαναλαμβάνεται χωρίς
αυτήν. Η διαδικασία επαναλαμβάνεται έως ότου ο έλεγχος $\chi^2$ πάψει να ανιχνεύει εσφαλμένα
δεδομένα. Στην πράξη χρησιμοποιείται συχνά το κατώφλι $r_i^{N} > 3$, δηλαδή υπόλοιπο
μεγαλύτερο από τρεις τυπικές αποκλίσεις του.

\subsubsection*{Δύο θεμελιώδεις περιορισμοί}

Η μέθοδος είναι ισχυρή, υπόκειται όμως σε δύο όρια που πρέπει να είναι σαφή.

\textbf{Πρώτον, οι κρίσιμες μετρήσεις.} Μια μέτρηση χαρακτηρίζεται \emph{κρίσιμη} όταν η
αφαίρεσή της καθιστά το σύστημα μη παρατηρήσιμο. Τέτοιες μετρήσεις εμφανίζουν εξ ορισμού
μηδενικό υπόλοιπο, καθώς δεν υφίσταται άλλη πληροφορία ικανή να τις διαψεύσει και η
εκτίμηση αναγκάζεται να τις αναπαραγάγει επακριβώς. Συνεπώς, σφάλμα σε κρίσιμη μέτρηση
είναι \emph{μη ανιχνεύσιμο}, ανεξαρτήτως μεγέθους. Ο εντοπισμός των κρίσιμων μετρήσεων ενός
δικτύου αποτελεί επομένως μέρος του σχεδιασμού του συστήματος μέτρησης και όχι ακαδημαϊκή
άσκηση.

\textbf{Δεύτερον, οι επιθέσεις έγχυσης ψευδών δεδομένων.} Οι Liu, Ning και Reiter ανέδειξαν
το 2009 ένα ουσιωδώς σοβαρότερο ζήτημα. Έστω επιτιθέμενος με γνώση της τοπολογίας και των
παραμέτρων του δικτύου, δηλαδή του πίνακα $\mathbf{H}$. Εφόσον αλλοιώσει τις μετρήσεις κατά
$\mathbf{a} = \mathbf{H}\mathbf{c}$ για οποιοδήποτε διάνυσμα $\mathbf{c}$, η αλλοιωμένη
μέτρηση $\mathbf{z}_{\text{bad}} = \mathbf{z} + \mathbf{a}$ αντιστοιχεί επακριβώς στην
κατάσταση $\hat{\mathbf{x}} + \mathbf{c}$, ενώ το υπόλοιπο παραμένει \emph{αμετάβλητο}:

\begin{equation}
\mathbf{r}_{\text{bad}} = \mathbf{z} + \mathbf{H}\mathbf{c}
        - \mathbf{H}(\hat{\mathbf{x}} + \mathbf{c}) = \mathbf{r}
\label{eq:fdia}
\end{equation}

Η επίθεση διαφεύγει πλήρως των ανωτέρω στατιστικών ελέγχων, ακριβώς επειδή είναι
\emph{συνεπής} προς το μοντέλο. Δεν πρόκειται για ελάττωμα υλοποίησης, αλλά για δομικό όριο
κάθε ελέγχου που βασίζεται στα υπόλοιπα. Οι επιθέσεις αυτές, γνωστές ως επιθέσεις έγχυσης
ψευδών δεδομένων (False Data Injection Attacks, FDIA), καθώς και οι στρατηγικές
αντιμετώπισής τους, εξετάζονται αναλυτικά στο Κεφάλαιο 8.

\subsubsection*{Ενσωμάτωση μετρήσεων PMU}

Η προσθήκη μετρήσεων PMU μεταβάλλει ουσιωδώς το πρόβλημα. Καθώς η PMU μετρά απευθείας τη
γωνία τάσης, η σχέση μεταξύ μετρήσεων και κατάστασης καθίσταται \emph{γραμμική} όταν
εκφραστεί σε καρτεσιανές (ορθογώνιες) συντεταγμένες. Εφόσον το σύστημα είναι πλήρως παρατηρήσιμο μόνο
από PMU, η \eqref{eq:measmodel} εκφυλίζεται σε
$\mathbf{z} = \mathbf{H}\mathbf{x} + \mathbf{e}$ και η λύση δίνεται σε ένα βήμα:

\begin{equation}
\hat{\mathbf{x}} = \left(\mathbf{H}^{T}\mathbf{R}^{-1}\mathbf{H}\right)^{-1}
                    \mathbf{H}^{T}\mathbf{R}^{-1}\mathbf{z}
\label{eq:linse}
\end{equation}

Τα πλεονεκτήματα είναι σημαντικά: δεν απαιτούνται επαναλήψεις, οπότε δεν ανακύπτουν
προβλήματα σύγκλισης ούτε εξάρτηση από την αρχική εκτίμηση· ο $\mathbf{H}$ παραμένει
σταθερός για δεδομένη τοπολογία και συνεπώς προϋπολογίζεται· τέλος, η εκτίμηση μπορεί να
εκτελείται με τον ρυθμό των PMU, δηλαδή δεκάδες φορές ανά δευτερόλεπτο αντί για μία φορά
ανά λεπτό.

Στην πράξη, η πλήρης κάλυψη με PMU σπανίζει για λόγους κόστους, οπότε χρησιμοποιούνται
\textbf{υβριδικοί εκτιμητές} που συνδυάζουν συμβατικές μετρήσεις SCADA με
συγχρονισμένους φασιθέτες. Η βασική δυσκολία έγκειται στην ετερογένεια των ρυθμών, δηλαδή στον
συνδυασμό μετρήσεων που ανανεώνονται ανά 4 s με άλλες που ανανεώνονται ανά 20 ms.

\subsection{Στάδια λειτουργίας: από τον προγραμματισμό στον πραγματικό χρόνο}

Η λειτουργία του συστήματος δεν αποφασίζεται σε μία χρονική στιγμή, αλλά μέσα από
διαδοχικές χρονικές κλίμακες. Η λογική είναι απλή και αξίζει να διατυπωθεί ρητά: \emph{όσο
πλησιάζει η ώρα παράδοσης, τόσο ακριβέστερη καθίσταται η πρόβλεψη και τόσο δαπανηρότερη η
διόρθωση}. Κάθε στάδιο διορθώνει τα σφάλματα του προηγούμενου, με λιγότερο διαθέσιμο χρόνο
και περιορισμένες επιλογές. Στο ευρωπαϊκό Μοντέλο-Στόχο (Target Model) οι κλίμακες αυτές
αντιστοιχούν σε διακριτές αγορές.

\begin{itemize}
\item \textbf{Μακροπρόθεσμος και προθεσμιακός ορίζοντας} (έτη έως εβδομάδες). Μελέτες
      επάρκειας ισχύος, συντονισμός προγραμμάτων συντήρησης γραμμών και μονάδων, διαστασιολόγηση
      των εφεδρειών, συναλλαγές στην Προθεσμιακή Αγορά για αντιστάθμιση του κινδύνου τιμής.
\item \textbf{Αγορά Επόμενης Ημέρας} (day-ahead). Αποτελεί το κύριο στάδιο
      προγραμματισμού. Οι συμμετέχοντες υποβάλλουν εντολές αγοράς και πώλησης για κάθε
      αγοραία χρονική μονάδα της επόμενης ημέρας και η αγορά εκκαθαρίζεται μέσω της ενιαίας
      σύζευξης επόμενης ημέρας σε ευρωπαϊκό επίπεδο (Single Day-Ahead Coupling, SDAC),
      ώστε η ενέργεια να ρέει από τις φθηνότερες προς τις ακριβότερες ζώνες προσφοράς έως
      την εξάντληση της διαθέσιμης διαζωνικής δυναμικότητας. Ο
      Διαχειριστής ελέγχει στη συνέχεια την τεχνική εφικτότητα του προκύπτοντος
      προγράμματος.
\item \textbf{Ενδοημερήσια Αγορά} (intra-day). Επιτρέπει στους συμμετέχοντες τη διόρθωση των θέσεών τους καθώς
      πλησιάζει η ώρα παράδοσης και βελτιώνονται οι προβλέψεις. Έχει ιδιαίτερη σημασία σε
      συστήματα με υψηλή διείσδυση αιολικής και φωτοβολταϊκής παραγωγής: το σφάλμα
      πρόβλεψης αιολικής παραγωγής μειώνεται δραστικά όσο συντομεύει ο ορίζοντας, οπότε η
      δυνατότητα αναπροσαρμογής λίγες ώρες πριν από την παράδοση αποκτά σημαντική
      οικονομική αξία.
\item \textbf{Πραγματικός χρόνος και εξισορρόπηση} (balancing). Ο Διαχειριστής ενεργοποιεί
      προσφορές ενέργειας εξισορρόπησης, ώστε να καλύψει τις εναπομένουσες
      ανισορροπίες, να
      αντιμετωπίσει διαταραχές και να διαχειριστεί συμφορήσεις.
\end{itemize}

\subsubsection*{Η αγοραία χρονική μονάδα}

Η \emph{αγοραία χρονική μονάδα} (Market Time Unit, MTU) είναι η χρονική ανάλυση με την
οποία εκκαθαρίζεται η αγορά. Επί σειρά ετών ήταν η ώρα. Από τις 30 Σεπτεμβρίου 2025, με ημέρα παράδοσης την 1η
Οκτωβρίου 2025, η ενιαία σύζευξη επόμενης ημέρας λειτουργεί με αγοραία χρονική μονάδα
\textbf{15 λεπτών} σε όλες τις ευρωπαϊκές ζώνες προσφοράς, εναρμονισμένη με την περίοδο
εκκαθάρισης ανισορροπιών. Η μεταβολή δεν συνιστά τεχνική λεπτομέρεια: μικρότερη αγοραία
χρονική μονάδα συνεπάγεται ότι οι ταχείες μεταβολές παραγωγής και φορτίου αποτυπώνονται στο ίδιο
το πρόγραμμα, αντί να μετακυλίονται στην εξισορρόπηση, με αποτέλεσμα τη μείωση των αναγκών
σε εφεδρείες.

\subsubsection*{Το ελληνικό πλαίσιο}

Στην Ελλάδα το Μοντέλο-Στόχος τέθηκε σε εφαρμογή την 1η Νοεμβρίου 2020. Τη διαχείριση της
Αγοράς Επόμενης Ημέρας και της Ενδοημερήσιας Αγοράς έχει το Ελληνικό Χρηματιστήριο Ενέργειας
(ΕΧΕ), ενώ τη διαχείριση της \textbf{Αγοράς Εξισορρόπησης} ο ΑΔΜΗΕ. Ο τελευταίος εκτελεί τη Διαδικασία
Ολοκληρωμένου Προγραμματισμού και ενεργοποιεί σε πραγματικό χρόνο τόσο ενέργεια
εξισορρόπησης όσο και τις εφεδρείες που περιγράφονται στη συνέχεια.

\subsection{Απόκριση συχνότητας και εφεδρείες}

\subsubsection*{Αδράνεια και αρχικός ρυθμός μεταβολής}

Η συχνότητα αποτελεί τον πλέον άμεσο δείκτη του ισοζυγίου ισχύος: παραμένει σταθερή όσο
παραγωγή και ζήτηση εξισορροπούνται και αποκλίνει μόλις εμφανιστεί ανισοζύγιο.

Αμέσως μετά από ανισορροπία ισχύος $\Delta P$, τυπικά την απώλεια μεγάλης μονάδας
ηλεκτροπαραγωγής, κανένας
ρυθμιστής δεν έχει προλάβει να αντιδράσει. Το έλλειμμα καλύπτεται από την κινητική ενέργεια
των στρεφόμενων μαζών, οι οποίες επιβραδύνονται. Ο αρχικός ρυθμός μεταβολής της συχνότητας
(Rate of Change of Frequency, RoCoF) προκύπτει από την εξίσωση ταλάντωσης (swing equation)
των σύγχρονων μηχανών και δίνεται από τη σχέση

\begin{equation}
\frac{\mathrm{d}f}{\mathrm{d}t} = \frac{f_0 \, \Delta P}{2 H S_n}
\label{eq:rocof}
\end{equation}

όπου $f_0$ η ονομαστική συχνότητα, $H$ η σταθερά αδράνειας του συστήματος σε δευτερόλεπτα
και $S_n$ η συνολική ονομαστική φαινόμενη ισχύς των \emph{συγχρονισμένων} μονάδων.

\emph{Παράδειγμα.} Για συγχρονισμένη περιοχή με $S_n = 300$ GVA, $H = 5$ s και απώλεια
$\Delta P = 3\,000$ MW, ο αρχικός ρυθμός ανέρχεται σε
$50 \times 3000 / (2 \times 5 \times 300\,000) = 0{,}05$ Hz/s. Εφόσον η ίδια απώλεια
εκδηλωθεί σε ώρα χαμηλού φορτίου, με τη μισή συγχρονισμένη ισχύ, ο ρυθμός διπλασιάζεται.

Η σχέση \eqref{eq:rocof} εξηγεί τον λόγο για τον οποίο η αντικατάσταση σύγχρονων
γεννητριών από μονάδες που συνδέονται μέσω ηλεκτρονικών μετατροπέων ισχύος, όπως τα
αιολικά και τα φωτοβολταϊκά πάρκα — οι οποίες δεν προσφέρουν εγγενή αδράνεια — οδηγεί σε ταχύτερες αποκλίσεις συχνότητας. Ο διαθέσιμος χρόνος έως ότου η
συχνότητα φτάσει σε επικίνδυνα επίπεδα συρρικνώνεται και η απόκριση οφείλει να καταστεί
αντιστοίχως ταχύτερη. Από εδώ πηγάζει το ενδιαφέρον για συνθετική αδράνεια και ταχεία
απόκριση συχνότητας.

Το Σχήμα~\ref{fig:freq} δείχνει τη συνολική απόκριση για το συμβάν αναφοράς και
διακρίνει τις τρεις διαδοχικές φάσεις που αναλύονται στη συνέχεια.

\begin{figure}[!htbp]
\centering
\includegraphics[width=\linewidth]{figures/frequency-response.svg}
\caption{Απόκριση συχνότητας σε απώλεια παραγωγής 3\,000 MW. Η κλίση στην αρχή καθορίζεται
αποκλειστικά από την αδράνεια· η εφεδρεία διατήρησης συχνότητας (ΕΔΣ) ανακόπτει την πτώση
και σταθεροποιεί τη συχνότητα σε νέα μόνιμη τιμή· η εφεδρεία αποκατάστασης συχνότητας (ΕΑΣ)
εξαλείφει την εναπομένουσα
απόκλιση και επαναφέρει τα 50 Hz.}
\label{fig:freq}
\end{figure}

\subsubsection*{Εφεδρεία διατήρησης συχνότητας και στατισμός}

Η \textbf{εφεδρεία διατήρησης συχνότητας} (ΕΔΣ· Frequency Containment Reserve, FCR)
υλοποιείται από τους ρυθμιστές στροφών των μονάδων, οι οποίοι μεταβάλλουν την παραγόμενη
ισχύ αναλογικά προς την απόκλιση συχνότητας. Ο συντελεστής αναλογίας εκφράζεται μέσω του \emph{στατισμού} $R$:

\begin{equation}
\frac{\Delta P}{P_n} = -\frac{1}{R}\,\frac{\Delta f}{f_0}
\label{eq:droop}
\end{equation}

Τυπικές τιμές είναι $R = 4$ έως 5\%. \emph{Παράδειγμα:} μονάδα 500 MW με $R = 5\%$, υπό
απόκλιση $\Delta f = -200$ mHz, δηλαδή $-0{,}4\%$, αυξάνει την παραγωγή της κατά
$(0{,}004/0{,}05) \times 500 = 40$ MW.

Το αρνητικό πρόσημο εξασφαλίζει αρνητική ανάδραση, καθώς πτώση της συχνότητας προκαλεί
αύξηση της παραγωγής. Ουσιώδες είναι ωστόσο ότι η ρύθμιση είναι \emph{αναλογική} και όχι
ολοκληρωτική: για να διατηρηθεί η επιπλέον παραγωγή, οφείλει να διατηρηθεί και η απόκλιση
συχνότητας. Η ΕΔΣ επομένως \textbf{σταθεροποιεί} τη συχνότητα σε νέα μόνιμη
τιμή, χωρίς να την \textbf{επαναφέρει} στα 50 Hz. Ακριβώς αυτό το μόνιμο σφάλμα
καλείται να εξαλείψει το επόμενο επίπεδο.

\subsubsection*{Εφεδρεία αποκατάστασης συχνότητας σε απομονωμένο σύστημα}

Η \textbf{εφεδρεία αποκατάστασης συχνότητας} (ΕΑΣ· Frequency Restoration Reserve, FRR) εξετάζεται
πρώτα στην απλούστερη περίπτωση: ένα \emph{απομονωμένο} σύστημα, χωρίς διασυνδέσεις με
γειτονικά συστήματα.

Εκεί υπάρχει ένα και μόνο κριτήριο, η συχνότητα. Ο κεντρικός ελεγκτής εφαρμόζει
\emph{ολοκληρωτική} δράση επί της απόκλισης:

\begin{equation}
\Delta P_{\mathrm{FRR}}(t) = -K_I \int_0^{t} \Delta f(\tau)\,\mathrm{d}\tau
\label{eq:frr-isolated}
\end{equation}

Ο ολοκληρωτικός όρος είναι εκείνος που κάνει τη διαφορά. Όσο υπάρχει απόκλιση συχνότητας,
το ολοκλήρωμα συνεχίζει να αυξάνεται και η εντολή παραγωγής να μεγαλώνει· η διαδικασία
σταματά μόνο όταν $\Delta f = 0$. Με άλλα λόγια, η ολοκληρωτική δράση εξαλείφει εξ ορισμού
το μόνιμο σφάλμα που αφήνει η αναλογική δράση της ΕΔΣ. Καθώς η συχνότητα επανέρχεται στα
50 Hz, η συνεισφορά της ΕΔΣ μηδενίζεται σύμφωνα με την \eqref{eq:droop}: η ΕΑΣ
\emph{αποδεσμεύει} την ΕΔΣ, ώστε αυτή να είναι
και πάλι διαθέσιμη για την επόμενη διαταραχή.

Η δράση αυτή είναι σκόπιμα βραδύτερη από την ΕΔΣ. Δύο ελεγκτές που δρουν στο ίδιο μέγεθος με
συγκρίσιμες ταχύτητες αλληλεπιδρούν και μπορούν να οδηγήσουν σε ταλαντώσεις· ο διαχωρισμός
των χρονικών κλιμάκων — δευτερόλεπτα για την ΕΔΣ, λεπτά για την ΕΑΣ — είναι εκείνος που
εξασφαλίζει την ευστάθεια της συνολικής ιεραρχίας.

\subsubsection*{Εφεδρεία αποκατάστασης συχνότητας σε σύστημα πολλών περιοχών}

Σε διασυνδεδεμένο σύστημα το πρόβλημα αλλάζει ποιοτικά. Η συχνότητα είναι \emph{κοινή} σε
ολόκληρη τη συγχρονισμένη περιοχή: μια απώλεια μονάδας στη Γερμανία μειώνει τη συχνότητα και
στην Ελλάδα. Αν κάθε Διαχειριστής ρύθμιζε μόνο με βάση τη συχνότητα, όλοι θα αντιδρούσαν
ταυτόχρονα στην ίδια διαταραχή, ανεξαρτήτως του πού συνέβη. Το αποτέλεσμα θα ήταν
υπερδιόρθωση, ανεξέλεγκτες αποκλίσεις των διασυνοριακών ροών από τα προγράμματα και εν τέλει
οικονομική μεταφορά κόστους από τον υπαίτιο στους υπόλοιπους.

Απαιτείται επομένως δεύτερο κριτήριο, το οποίο να διακρίνει \emph{πού} εκδηλώθηκε το
ανισοζύγιο. Ο ελεγκτής κάθε περιοχής ελέγχου φορτίου-συχνότητας (περιοχή ΕΦΣ· LFC area) επιδιώκει τον
μηδενισμό του \emph{σφάλματος ελέγχου περιοχής} (ΣΕΠ· Area Control Error, ACE):

\begin{equation}
\mathrm{ACE} = \Delta P_{\mathrm{tie}} + K \, \Delta f
\label{eq:ace}
\end{equation}

όπου $\Delta P_{\mathrm{tie}}$ το σφάλμα ελέγχου ισχύος, δηλαδή η απόκλιση της μετρούμενης
ανταλλαγής ισχύος στις γραμμές διασύνδεσης από το πρόγραμμα, και $K$ ο συντελεστής K της
περιοχής σε MW/Hz (στην κλασική βιβλιογραφία συντελεστής πόλωσης συχνότητας, frequency
bias). Στο απομονωμένο σύστημα ο πρώτος όρος δεν υπάρχει και η \eqref{eq:ace} εκφυλίζεται
στο κριτήριο της προηγούμενης υποενότητας. Για τον λόγο αυτόν ο SO GL ορίζει ενιαία ως
μέγεθος ελέγχου της αποκατάστασης το \emph{σφάλμα ελέγχου αποκατάστασης συχνότητας} (ΣΕΑΣ·
Frequency Restoration Control Error, FRCE): σε σύστημα πολλών περιοχών ισούται με το ΣΕΠ
της περιοχής, ενώ όταν η περιοχή ελέγχου συμπίπτει με ολόκληρη τη συγχρονισμένη περιοχή
ισούται με την απόκλιση συχνότητας.

Ο συνδυασμός των δύο όρων παράγει την επιθυμητή διάκριση, όπως δείχνει το
Σχήμα~\ref{fig:lfc}:

\begin{itemize}
\item \textbf{Διαταραχή εντός της περιοχής.} Η συχνότητα πέφτει ($\Delta f < 0$) και,
      ταυτόχρονα, οι γειτονικές περιοχές παρέχουν ισχύ υποστήριξης, οπότε οι εισαγωγές υπερβαίνουν το
      πρόγραμμα ($\Delta P_{\mathrm{tie}} < 0$). Οι δύο όροι έχουν το ίδιο πρόσημο, το ΣΕΠ
      λαμβάνει μεγάλη απόλυτη τιμή και η περιοχή διορθώνει.
\item \textbf{Διαταραχή εκτός της περιοχής.} Η συχνότητα πέφτει επίσης, αλλά η περιοχή
      \emph{εξάγει} επιπλέον ισχύ προς τη γειτονική περιοχή ($\Delta P_{\mathrm{tie}} > 0$). Με
      κατάλληλη επιλογή του $K$ οι δύο όροι αλληλοαναιρούνται και το ΣΕΠ παραμένει σχεδόν
      μηδενικό: η περιοχή έχει ήδη προσφέρει τη συνεισφορά της μέσω της ΕΔΣ και δεν
      απαιτείται περαιτέρω ενέργεια.
\end{itemize}

Η ιδιότητα αυτή ονομάζεται \emph{μη αλληλεπιδραστικός έλεγχος} (non-interactive control):
κάθε Διαχειριστής αποκαθιστά τα δικά του ανισοζύγια, ενώ η αλληλεγγύη κατά τα πρώτα
δευτερόλεπτα της διαταραχής παραμένει καθολική. Ο συντελεστής $K$ επιλέγεται όσο το δυνατόν
πλησιέστερα στη συνολική απόκριση ισχύος της περιοχής ανά μονάδα απόκλισης συχνότητας —
άθροισμα της συνεισφοράς της ΕΔΣ, του αυτοελέγχου της παραγωγής και της αυτορρύθμισης του
φορτίου — ώστε η αλληλοαναίρεση να είναι όσο το δυνατόν πληρέστερη.

\begin{figure}[!htbp]
\centering
\includegraphics[width=\linewidth]{figures/lfc-areas.svg}
\caption{Εφεδρεία αποκατάστασης συχνότητας σε απομονωμένο και σε διασυνδεδεμένο σύστημα. Στην πρώτη
περίπτωση αρκεί η απόκλιση συχνότητας· στη δεύτερη απαιτείται και ο όρος των ανταλλαγών ισχύος
στις γραμμές διασύνδεσης, ώστε να διορθώνει μόνο η περιοχή στην οποία εκδηλώθηκε η διαταραχή.}
\label{fig:lfc}
\end{figure}

\subsubsection*{Τα τέσσερα επίπεδα εφεδρείας}

\begin{itemize}
\item \textbf{ΕΔΣ} (FCR) --- εφεδρεία διατήρησης συχνότητας. Αυτόματη και τοπική,
      ενεργοποιείται από κοινού σε ολόκληρη τη συγχρονισμένη περιοχή. Συγκρατεί και
      σταθεροποιεί τη συχνότητα.
\item \textbf{αΕΑΣ} (aFRR) --- εφεδρεία αποκατάστασης συχνότητας με αυτόματη ενεργοποίηση.
      Κεντρική, οδηγείται από το ΣΕΑΣ μέσω της ΑΡΠ. Επαναφέρει τη συχνότητα στα 50 Hz και
      αποδεσμεύει την ΕΔΣ.
\item \textbf{χΕΑΣ} (mFRR) --- εφεδρεία αποκατάστασης συχνότητας με χειροκίνητη
      ενεργοποίηση. Ενεργοποιείται με εντολή κατανομής του Διαχειριστή και αποδεσμεύει την
      αΕΑΣ σε παρατεταμένα ανισοζύγια.
\item \textbf{ΕΑ} (RR) --- εφεδρεία αντικατάστασης. Αποκαθιστά το απαιτούμενο επίπεδο των
      ενεργοποιημένων ΕΑΣ, ώστε το σύστημα να είναι έτοιμο για την επόμενη διαταραχή.
\end{itemize}

Στην παλαιότερη βιβλιογραφία και σε μέρος της πρακτικής τα τρία πρώτα επίπεδα απαντούν ως
πρωτεύουσα, δευτερεύουσα και τριτεύουσα ρύθμιση αντιστοίχως. Η ορολογία που χρησιμοποιείται
εδώ είναι εκείνη του SO GL, με τα ακρωνύμια της επίσημης ελληνικής απόδοσής του που
χρησιμοποιεί και ο ΑΔΜΗΕ, και προτιμάται, καθώς ονομάζει τα επίπεδα κατά τη \emph{λειτουργία}
τους — διατήρηση, αποκατάσταση, αντικατάσταση — και όχι κατά τη σειρά ενεργοποίησης.

Διακρίνεται ένα σαφές μοτίβο: κάθε επίπεδο είναι βραδύτερο, οικονομικότερο και πιο
στοχευμένο από το προηγούμενο, ενώ ο ρόλος του συνίσταται στην \emph{αποδέσμευση} του
ταχύτερου, ώστε αυτό να καταστεί εκ νέου διαθέσιμο.

\begin{table}[!htbp]
\centering
\begin{tabular}{lllc}
\hline
Εφεδρεία & Λειτουργία & Ενεργοποίηση & Χρόνος πλήρους ενεργοποίησης \\
\hline
ΕΔΣ (FCR)   & διατήρηση      & αυτόματη, τοπική   & 30 s \\
αΕΑΣ (aFRR) & αποκατάσταση   & αυτόματη, κεντρική & 5 min \\
χΕΑΣ (mFRR) & αποκατάσταση   & χειροκίνητη        & 12,5 min \\
ΕΑ (RR)     & αντικατάσταση  & χειροκίνητη        & $>$ 15 min \\
\hline
\end{tabular}
\caption{Επίπεδα εφεδρείας συχνότητας κατά τον SO GL. Οι χρόνοι πλήρους ενεργοποίησης
αφορούν τη συγχρονισμένη περιοχή της Ηπειρωτικής Ευρώπης· για την αΕΑΣ και τη χΕΑΣ είναι
εκείνοι των τυποποιημένων προϊόντων των ευρωπαϊκών πλατφορμών εξισορρόπησης.}
\label{tab:reserves}
\end{table}

\subsubsection*{Ποσοτικές απαιτήσεις για την Ηπειρωτική Ευρώπη}

\begin{itemize}
\item Το \textbf{συμβάν αναφοράς} (reference incident) για τη διαστασιολόγηση της ΕΔΣ είναι
      ανισορροπία 3\,000 MW, είτε ως απώλεια παραγωγής είτε ως απώλεια φορτίου. Η ελάχιστη
      διαθέσιμη ΕΔΣ στη συγχρονισμένη περιοχή ανέρχεται ανά πάσα στιγμή σε 3\,000 MW.
\item Η \textbf{μέγιστη απόκλιση συχνότητας σταθερής κατάστασης} για το συμβάν αναφοράς
      είναι 200 mHz. Από τα δύο αυτά μεγέθη προκύπτει ο συντελεστής αναλογίας μεταξύ ΕΔΣ
      και απόκλισης συχνότητας: $3\,000\ \text{MW} / 0{,}2\ \text{Hz} = 15\,000$ MW/Hz.
\item Η \textbf{τυπική περιοχή συχνότητας} (standard frequency range) είναι $\pm 50$ mHz·
      η συχνότητα δεν επιτρέπεται να βρίσκεται εκτός αυτής για περισσότερα από
      15\,000 λεπτά ανά έτος. Η μέγιστη στιγμιαία απόκλιση
      ορίζεται σε 800 mHz.
\item Απόκλιση συχνότητας σταθερής κατάστασης μεγαλύτερη των 200 mHz κατατάσσει το
      σύστημα σε κατάσταση \emph{έκτακτης ανάγκης}.
\end{itemize}

\subsection{Σύνοψη}

Η λειτουργία του ΣΜΗΕ στηρίζεται στον κλειστό βρόχο με τον οποίο ξεκίνησε το κεφάλαιο:

\begin{enumerate}
\item Η \textbf{υποδομή εποπτείας} (SCADA, PMU) παράγει μετρήσεις, οι οποίες είναι
      πλεονάζουσες, θορυβώδεις και εν μέρει ασύγχρονες.
\item Η \textbf{εκτίμηση κατάστασης} τις μετατρέπει σε συνεπή και στατιστικά βέλτιστη
      εικόνα του δικτύου, απορρίπτοντας ταυτόχρονα τα προφανώς εσφαλμένα δεδομένα.
\item Τα \textbf{εργαλεία ανάλυσης επιχειρησιακής ασφάλειας} αξιολογούν την εικόνα αυτή
      έναντι του κριτηρίου N-1 και κατατάσσουν το σύστημα σε μία από τις πέντε καταστάσεις
      του συστήματος.
\item Οι \textbf{αποφάσεις} υλοποιούνται μέσω των αγορών και των διατάξεων ελέγχου, σε
      κλίμακες που εκτείνονται από τα δευτερόλεπτα της ΕΔΣ έως τους μήνες
      του προγραμματισμού.
\end{enumerate}

Ο βρόχος είναι τόσο ισχυρός όσο ο ασθενέστερος κρίκος του, και ο κρίκος αυτός είναι συχνά
η αξιοπιστία της ίδιας της εικόνας: κρίσιμες μετρήσεις χωρίς πλεονασμό, επιθέσεις συνεπείς
προς το μοντέλο, τμήματα του δικτύου που δεν είναι παρατηρήσιμα.

Οι υπηρεσίες που στηρίζουν τον βρόχο αυτόν ονομάζονται \emph{επικουρικές υπηρεσίες} και
διακρίνονται σε σχετικές και μη σχετικές με τη συχνότητα· τις εξετάζουν αντιστοίχως τα
Κεφάλαια 3 και 2. Το αντίστοιχο πρόβλημα στο επίπεδο της διανομής, όπου η παρατηρησιμότητα
είναι πολύ χαμηλότερη και κυριαρχούν οι ψευδομετρήσεις, εξετάζεται στο Κεφάλαιο 4. Η
ασφάλεια της ίδιας της υποδομής μέτρησης και ελέγχου αποτελεί το αντικείμενο του
Κεφαλαίου 8.
